%% pc-securite-abaque — effets du courant électrique sur le corps humain.
%% Abaque à échelles logarithmiques (durée en ordonnée, intensité en abscisse) :
%% les graduations ne sont PAS régulières, leurs positions sont posées à la main
%% dans les deux listes \foreach ci-dessous (une position en cm par valeur).
%% Cinq zones colorées ; les frontières des trois dernières sont incurvées
%% (elles se décalent vers la droite quand la durée diminue).
\documentclass[tikz,border=4pt]{standalone}
\usepackage{liaisons}
\begin{document}
%% --- couleurs des cinq zones (du moins au plus grave) -----------------------
\definecolor{zvert}{RGB}{30,150,90}
\definecolor{zjaune}{RGB}{245,205,40}
\definecolor{zorange}{RGB}{235,150,20}
\definecolor{zorangef}{RGB}{225,105,20}
\definecolor{zrouge}{RGB}{200,30,25}

%% l'unité x vaut 0,86 cm : le graphe reste sous 8 cm de large sans changer les textes
\begin{tikzpicture}[x=0.86cm,y=1cm,
  axe/.style={-{Stealth[length=5pt]}, line width=0.7pt},
  grad/.style={font=\tiny},
  grille/.style={draw=white, line width=0.35pt},
  %% noms des zones : écrits verticalement, sur deux lignes quand la zone est étroite
  %% (ils doivent rester dans la partie haute, où les frontières sont verticales)
  nomzone/.style={font=\tiny\bfseries, rotate=90, anchor=east, align=center},
  legende/.style={font=\scriptsize, anchor=west, inner sep=0pt}]

\def\XM{7.2}   % largeur du cadre
\def\YM{3.40}  % hauteur du cadre

%% ================= zones (empilées de la plus grave à la moins grave) =======
%% Chaque zone est peinte sur toute la largeur restante : il suffit alors de
%% décrire la frontière GAUCHE de la zone suivante, du bas vers le haut.
\fill[zrouge] (0,0) rectangle (\XM,\YM);
\fill[zorangef] (0,0) -- (6.40,0) .. controls (5.85,0.75) and (5.30,1.60) ..
      (5.30,\YM) -- (0,\YM) -- cycle;
\fill[zorange] (0,0) -- (4.80,0) .. controls (4.15,0.75) and (3.60,1.60) ..
      (3.60,\YM) -- (0,\YM) -- cycle;
\fill[zjaune] (0,0) -- (3.90,0) .. controls (3.35,0.75) and (3.00,1.60) ..
      (3.00,\YM) -- (0,\YM) -- cycle;
\fill[zvert] (0,0) rectangle (1.50,\YM);

%% ================= grille blanche fine =====================================
\foreach \gx in {0.55,1.05,1.50,1.95,2.50,3.00,3.60,4.40,5.30,6.10} {
  \draw[grille] (\gx,0) -- (\gx,\YM); }
\foreach \gy in {0.45,0.82,1.13,1.45,1.85,2.18,2.50,2.90} {
  \draw[grille] (0,\gy) -- (\XM,\gy); }

%% ================= noms des zones (en haut, écrits verticalement) ==========
\node[nomzone, text=white]  at (0.75,\YM-0.14) {Picotement};
\node[nomzone, text=black]  at (2.28,\YM-0.14) {Tétanisation};
\node[nomzone, text=white]  at (3.30,\YM-0.14) {Paralysie\\respiratoire};
\node[nomzone, text=white]  at (4.45,\YM-0.14) {Fibrillation\\ventriculaire};
\node[nomzone, text=white]  at (6.30,\YM-0.14) {Arrêt cardiaque};

%% ================= repère blanc : 30 mA pendant 500 ms =====================
\draw[black!55, line width=2.1pt] (0,1.85) -- (3.60,1.85) -- (3.60,0);
\draw[white, line width=1.1pt] (0,1.85) -- (3.60,1.85) -- (3.60,0);
\draw[white, line width=1.1pt, fill=black!55] (3.60,1.85) circle (0.085);

%% ================= cadre et axes ===========================================
\draw[line width=0.7pt, black!70] (0,0) rectangle (\XM,\YM);
\draw[axe] (0,0) -- (\XM+0.35,0);
\draw[axe] (0,0) -- (0,\YM+0.30);

%% graduations de l'axe des intensités (mA)
\foreach \gx/\lb in {0/0, 0.55/{0,2}, 1.05/{0,5}, 1.50/1, 1.95/2, 2.50/5,
                     3.00/10, 3.60/30, 4.40/100, 5.30/500, 6.10/{2\,000},
                     7.00/{10\,000}} {
  \draw[line width=0.5pt] (\gx,0) -- (\gx,-0.09);
  \node[grad, anchor=north, inner sep=2pt] at (\gx,-0.09) {\lb}; }
\node[font=\tiny, anchor=north] at ({0.5*\XM},-0.48)
  {Intensité du courant traversant le corps humain (mA)};

%% graduations de l'axe des durées (ms)
\foreach \gy/\lb in {0/0, 0.45/20, 0.82/50, 1.13/100, 1.45/200, 1.85/500,
                     2.18/{1\,000}, 2.50/{2\,000}, 2.90/{5\,000}, 3.25/{10\,000}} {
  \draw[line width=0.5pt] (0,\gy) -- (-0.09,\gy);
  \node[grad, anchor=east, inner sep=2pt] at (-0.09,\gy) {\lb}; }
\node[font=\tiny, rotate=90, align=center, anchor=south] at (-1.10,{0.5*\YM})
  {Durée de passage\\du courant (ms)};

%% ================= légende sous le graphe ==================================
\foreach \k/\co/\val/\txt in {%
    0/zrouge/{1 A}/{Arrêt du cœur},
    1/zorangef/{50/75 mA}/{Seuil de fibrillation cardiaque irréversible},
    2/zorange/{30 mA}/{Seuil de paralysie respiratoire au-delà de 500 ms},
    3/zjaune/{10 mA}/{Seuil de non-lâcher, contraction musculaire},
    4/zvert/{0,5 mA}/{Seuil de perception, sensation très faible}} {
  \pgfmathsetmacro\ly{-1.02-0.34*\k}
  \fill[\co] (-0.95,\ly) circle (0.095);
  \node[legende, font=\scriptsize\bfseries] at (-0.78,\ly) {\val};
  \node[legende] at (0.80,\ly) {\txt}; }
\end{tikzpicture}
\end{document}
